% sections/discussion.tex
\subsection*{Answers to the research questions}
Under this protocol, neither module improved thin-vessel segmentation (question 1); fine-tuned SAM gave no thin-vessel
advantage beyond training-seed variability and lost Dice and boundary accuracy (question 2); and Dice did not reveal the
thin-vessel differences, although it detected other consistent differences (question 3). Consistency across three seeds is the strongest statement three runs allow (one-sided $p = 0.05$), so
``no consistent effect'' means that this design could not establish an effect, not that none exists. Image-level
intervals, by contrast, exclude zero for differences of a few thousandths because they ignore seed variability; in post
hoc analyses, also resampling seeds widened them several-fold for thin-vessel recall.

\subsection*{Attention gates and multi-scale blocks}
The attention gates appear to have shifted the operating point rather than improved the segmentation: they raised recall
and lowered precision at the selected thresholds, did the reverse at 0.5, and thin-vessel recall followed; at matched
precision the attention effect was \PHStAttThinMatch. It also depended on the multi-scale blocks:
at 0.5, attention lowered thin-vessel recall about three times as much with them as without, so the main effect alone is
not a sufficient summary. The multi-scale decrease was consistent in direction, persisted at matched overall precision
and came with higher thin-vessel precision, so the blocks traded thin-vessel recall for precision rather than only
shifting the operating point. A design-specific explanation is plausible: the block splits its output channels between
three dilations, so only a third of the filters in each block keep the plain 3 $\times$ 3 receptive field that resolves
one-pixel structures. We tested neither this nor whether the gates were active; the identical stopping epochs of some same-seed model pairs (\nameref{si:runs}) suggest that shared
initialization and data order coupled their training. The results agree with the FIVES benchmark~\cite{fadugba2024benchmark} but apply to one design of
each module at 512 $\times$ 512 with soft targets and one learning rate; they do not show that attention or multi-scale
processing cannot help thin vessels.

\subsection*{Fine-tuned SAM and resolution}
Tiling changes several things at once: the detail in the input, the output grid, the part of the eye seen in one pass,
the target (binary tiles rather than soft masks) and the number of optimizer steps per epoch, and our design cannot
separate them. The mask-resampling references point to the output grid as a major contributor to the consistent
difference between tiled and whole-image SAM: the 256 $\times$ 256 grid alone kept \RefTwoThin{} of the thin-vessel
centreline and the 512 $\times$ 512 grid \RefFiveThin. This is a supported hypothesis, not a tested mechanism, and the
references are not upper bounds: whole-image SAM exceeded its 256 reference.

The comparison of tiled SAM with the CNNs is confounded in the same way. Its small gain in thin-vessel recall shrank to
\PHStTiledThinMatch{} at matched precision, so it mainly reflected a lower operating precision. Input resolution may
also contribute: tiled SAM sees the image at full resolution and the CNNs at 512, and in our one-seed validation
pilot a U-Net trained at 1024 gained \PilotThinGain{} over the 512 U-Net and almost reached the thin-vessel recall of the
seed-0 tiled SAM (\PilotTenThin{} against \SAMTiledValThinSeedA, both at threshold \PilotThreshold) at higher precision
(\PilotTenPrec{} against \PilotSAMTiledPrec). With one seed on validation data this is suggestive only;
separating input resolution and binary targets from the foundation model needs a resolution-matched control with several
seeds. Training may also have limited SAM: most runs reached the 150-epoch limit, with an
untuned learning rate. Zero-shot SAM failed even with
ground-truth points, as reported for vessels before~\cite{shi2023samzeroshot,qiu2023learnable}.

\subsection*{What overlap metrics show}
Dice detected consistent differences, for example between every SAM configuration and the U-Net, but ranked the SAM
configurations and the disease categories differently from thin-vessel recall. This is expected of an overlap score when
thin vessels hold few pixels~\cite{reinke2024pitfalls}. Dice also governed checkpoint and threshold
selection: runs with almost the same best validation Dice differed in thin-vessel recall, and selecting the threshold by
Dice reversed the sign of the attention effect. Whether thin-vessel
recall varied within a run while validation Dice did not is unknown, because we recorded neither thin-vessel recall nor
checkpoints at every epoch.

\subsection*{Why thin-vessel recall stays low}
All automatic models recovered only \AllThinMin--\AllThinMax{} of the thin-vessel centreline. Our data suggest, but do not
test, four contributing factors. First, thin-vessel labels may be uncertain: thin-vessel precision was also low for every
model, and where two masks exist for one image, each
recovered on average \AgreeThinMean{} of the other's thin-vessel centreline, although this evidence is post hoc, limited to
\AgreeNPairs{} AMD images. Second, for the CNNs and whole-image SAM the loss pushes soft
thin-vessel targets towards low probabilities (a target of 0.25 towards about \LossOptQuarter), consistent with the steep
dependence of thin-vessel recall on the threshold. Third, input resolution limits the CNNs, as the 1024 pilot suggests,
whereas the 512 output grid alone kept \RefFiveThin{} of the thin-vessel centreline. Fourth, selection by Dice does not
favour thin vessels.

\subsection*{Comparison with published results and recommendations}
Published Dice values on FIVES are not directly comparable with ours, because scoring region, resolution, split and
threshold rule differ; for example, Cao et al.\ report a Dice of 0.8404 for a multi-scale attention
U-Net~\cite{cao2023mfaunet}, below our U-Net's \UNetDiceMean, which may reflect the protocol rather than the model. For
comparisons of vessel segmentation methods, our results suggest reporting a width-stratified measure, with its resolution,
alongside overlap scores; training several seeds and reporting seed variability separately from image-level statistics;
reporting more than one threshold rule, the full threshold curve or comparisons at matched precision; and auditing the
dataset before training.

\subsection*{Limitations}
Three seeds estimate seed variability only roughly; the main image-level tests do not include it, and the post hoc
intervals that do rest on three runs per model. The test images are not independent: they include a duplicated image, a
shared mask and possible same-patient pairs, and \NIsoWithTest{} such pairs link the official training images to the test
set. All data come from one hospital and from Asian patients~\cite{jin2022fives}, and there was no external testing.
Several SAM runs reached the epoch limit with an untuned learning rate; SAM's second and third seeds were added after the
first test evaluation, and its comparisons are confounded by input resolution and targets, controlled only by a one-seed
pilot. The CNN results hold for one design of each module, 512
$\times$ 512 inputs and soft targets, which at a threshold of 0.5 keep thin vessels less well than the binary
alternatives (\nameref{si:resize}); we tested no other targets and did not inspect the attention coefficients.
The post hoc analyses were added after the main test evaluation. Metrics were computed inside an FOV mask, NSD is a boundary-pixel approximation, and vessel
density and fractal dimension are image-derived measures without clinical validation. The cost figures come from single
runs on one GPU type, and the annotation-agreement analysis is post hoc and small.

\subsection*{Future work}
Next steps follow from these limitations: a CNN trained at full resolution with several seeds, on whole images or on the
same tiles with binary targets, as a resolution-matched control for tiled SAM; longer SAM training with a tuned learning
rate; inspection of the attention coefficients; and training targets or losses that keep thin vessels. External testing
and a self-configuring baseline such as nnU-Net~\cite{isensee2021nnunet} would show how far the findings generalize.
